\section{Computational Complexity}\label{sec:complexity}
% \subsection{Computational Complexity}\label{sec:complexity}

Rather than measuring raw execution time, which depends heavily on hardware and software, we count the number of cryptographic operations as these are the computationally most expensive part of the scheme.
Table~\ref{tab:operation_count} summarizes this asymptotic counting for each party, considering a path of $m$ nodes and $t$ semi-trusted third parties (STTPs).  
The total cost of sending a single message can be obtained by multiplying the cost per STTP by $t$ and the cost per node by $m$.

The client bears most of the overhead introduced by STTPs, performing \bigO{m t}\footnote{\label{cnote}
    In our current implementation, random EC points for the shares are generated by multiplying the curve generator $G$ by a random scalar.  
    This could instead be replaced by mapping these random scalars to EC points directly using Elligator, which is faster than performing EC multiplications.
    Adopting this optimization would reduce the client's asymptotic cost for EC multiplications to \bigO{m},  
    but at the expense of increasing the asymptotic cost of the Elligator integer-to-point mapping to \bigO{m t}.
}
elliptic-curve multiplications and additions to split $m$ addresses and cryptographic secrets into $t$ shares.  
For each node, the client also computes one hash of the shared secret, one Elligator mapping to encode its IP, and one inverse mapping \eqref{eq:shared_secret}, all of which scale linearly with $m$.

For each packet layer, a STTP performs EC multiplications, additions, and hashes proportional to the number of chunks \eqref{eq:preprocessing}, \eqref{eq:layer_i},  and \eqref{eq:integrity}.  
As both the number of chunks and layers increase linearly with $m$, the total computational cost per STTP scales quadratically as $\mathcal{O}(p^2)$.  
However, since STTPs operate in parallel, their collective workload has limited impact on the overall latency.

Each node performs EC multiplications, additions and hashes proportional only to the number of chunks \eqref{eq:integrity_check} and \eqref{eq:decrypt}, yielding linear complexity in $m$.  
Additionally, each node executes one Elligator inverse mapping \eqref{eq:derive_secret} to convert the shared secret point into a shared secret integer.

In practice, most anonymous communication networks implementations such as onion routing~\cite{onion-routing96} and mixnets~\cite{nym} use routes of three to five nodes, making \bigO{p} $\approx$ \bigO{1}.
Accordingly, the overall computational burden is more sensitive to the number of STTPs $t$ than the path length $m$.
As our design remains secure as long as one STTP is honest, only a small number $t$ may be sufficient.
Determining the precise number needed to meet a desired trust threshold is left for future work.

\begin{table}[H]
\centering
\caption{Asymptotic number of cryptographic operations per party, with $t$ STTPs and a path of $m$ nodes.}
\label{tab:operation_count}
    \begin{tabular}{lccc}
        \toprule
        & \textbf{Client} & \textbf{STTP} & \textbf{Network Node} \\
        \midrule
        \textbf{EC Multiplication}              & \bigO{m \, t}\textsuperscript{\ref{cnote}}  & \bigO{m^{2}}  & \bigO{m} \\
        \textbf{EC Addition}                    & \bigO{m \, t}                             & \bigO{m^{2}}  & \bigO{m} \\
        \textbf{Elligator Mapping}              & \bigO{m}\textsuperscript{\ref{cnote}}       & \slash        & \slash \\
        \textbf{Inverse Mapping}                & \bigO{m}                                  & \slash        & \bigO{1} \\
        \textbf{Hash}                           & \bigO{m}                                  & \bigO{m^{2}}  & \bigO{m} \\
        \bottomrule
    \end{tabular}
\end{table} 

% \medskip
\noindent\textbf{Comparison with Sphinx.}
%
A direct comparison with the original Sphinx protocol is not entirely meaningful, as some primitives differ (e.g. Elligator).  
Nevertheless, from a performance perspective, DSphinx is expected to be slower than Sphinx.  
The main contributing factors are: (i) the extensive use of elliptic curve arithmetic, in contrast to the lightweight hashing and XOR operations in Sphinx; and (ii) the distributed nature of the header construction, which introduces communication latency.  
However, the computational load on individual hops in the paths remains comparable to Sphinx, and most of the added overhead is mitigated by the parallel execution of STTPs.  

Given that decentralized networks such as mixnets and Lightening inherently trade latency for privacy, the additional cost introduced by our decentralized design is reasonable.  
In return, it provides stronger flexibility for decentralized networks.
